% R4_FRONTMATTER, attempt 1: replacement front matter and Introduction only.
% Use with the supplied manuscript preamble and existing bibliography.
\begin{frontmatter}
\title{High-temperature property sensitivity after melt-pool length calibration in an IN625 conduction model}
% Authorship, affiliations and corresponding-author details require human confirmation.
\begin{abstract}
Melt-pool length provides a practical calibration target for laser-processing thermal models, but interpreting the fitted field requires distinguishing melting-boundary motion from the cooling experienced by material.
We examine this distinction with a nonlinear moving-frame enthalpy model for three bare-plate IN625 conditions from the Additive Manufacturing Metrology Testbed benchmark.
One effective source factor is fitted to surface length at condition B and retained throughout retrospective, nonblind geometry comparisons and a four-corner comparison of high-temperature conductivity and sensible heat-capacity continuations.
The length fit leaves the fusion envelope 19.54~\um{} wider and 4.88~\um{} shallower than the B reference means, and none of the tested continuations removes this discrepancy.
Relative to linear continuation of both properties, holding conductivity at its final tabulated value extends the pool by 43.21~\um{}, whereas holding sensible heat capacity shortens it by 6.17~\um{}, with latent heat retained.
The dominant length response is displacement of the rear solidus and liquidus together: rear solidus displacement accounts for 98.3\% of the conductivity-induced extension, while the separation of the rear boundaries changes much less.
At fixed power, increasing scan speed from 0.8 to 1.2~m/s produces rear phase spans near 80~\um{} but material passage times of 100.33 and 67.33~\us{}, respectively.
Thus sensitivity of total length chiefly reflects thermal-tail position in the tested property comparison, while cooling through the phase interval also depends on boundary separation and scan speed.
These model diagnostics distinguish changes in fusion shape, thermal-tail position and material cooling time; independent thermal measurements remain necessary to assess the physical cooling histories.
\end{abstract}
\begin{keyword}
Laser powder bed fusion \sep IN625 \sep Melt-pool geometry \sep Enthalpy model \sep High-temperature properties \sep Model discrepancy
\end{keyword}
\end{frontmatter}

\section{Introduction}
\label{sec:introduction}

% I1: The reader's problem: what a geometric fit says about cooling.
A thermal model can reproduce the length of a laser-melted track while leaving open how material cools behind the laser. This distinction matters when calculated temperature gradients and cooling histories supply the thermal conditions for a solidification or microstructure model. In laser powder bed fusion (LPBF), resolving laser coupling, phase change and liquid transport together is computationally demanding, motivating models with different levels of physical detail~\cite{king2015,review2024}. Melt-pool geometry supplies an accessible constraint on these calculations; thermography can add spatial temperature profiles and temporal information~\cite{hooper2018,myers2023}. The question is therefore what thermal information a geometric calibration actually constrains, and which features of the fitted field remain sensitive to material and source assumptions.

% I2: Model hierarchy and its output-dependent assessment.
The available model hierarchy offers several ways to balance cost and physical resolution. Thermo-fluid models resolve energy redistribution associated with surface-tension-driven flow, recoil pressure and evaporation~\cite{khairallah2016}. Moving-source conduction models instead describe heat transport without solving liquid momentum. Their lineage extends from Rosenthal's relations between source motion, heat flow and cooling to finite-source formulations with nonlinear enthalpy and phase change~\cite{rosenthal1946,vanelsen2007}. Scaling laws further simplify estimates of pool dimensions, with fidelity depending on material and operating regime~\cite{rubenchik2018,naderi2023}. Rapid transient conduction calculations have also been assessed against solidification conditions and grain structure in specific laser and electron-beam powder-bed cases~\cite{Plotkowski2017Rapid}. These results support affordable calculations for defined uses; their relevance to a cooling history depends on the observable and conditions used to assess them.

% I3: Geometric success through source and effective-transport choices.
Geometric prediction within this hierarchy depends on how the model represents both heating and transport. Analytical and finite-element models reasonably estimated IN718 melt-pool shapes, although their absorption sensitivities and responses at high heat input differed~\cite{promoppatum2017}. For IN625 in the Additive Manufacturing Metrology Testbed (AMMT) benchmark, measured laser input and calibrated effective directional conductivity improved the predicted dimensions~\cite{kollmannsberger2019}. Source calibration offers a complementary means of representing unresolved energy delivery: inverse fitting to boundary temperatures supports interpolation of double-ellipsoid parameters across process conditions~\cite{ross2022}, and a dynamic volumetric source calibrated over beam-size variation in bare-plate IN625 was subsequently applied to a layer-scale benchmark~\cite{dynamic2024}. Bayesian treatment of an adapted Eagar--Tsai model provides another route, propagating geometric uncertainty associated with apparent absorption and powder properties~\cite{whalen2021}. These approaches establish that useful geometry can be obtained through several source and transport descriptions. They also make the thermal interpretation of the chosen description a separate question from the success of its geometric fit.

% I4: A direct temperature precedent and observation/source definitions.
The need for an additional thermal constraint is demonstrated directly by the bare-plate 316L study of Myers et al.~\cite{myers2023}. Parameter combinations with similar agreement to cross-sectional geometry produced different surface-temperature profiles, and two-color imaging helped discriminate among them. The earlier AMMT model likewise explicitly limited the validity of its internal melt-pool temperatures despite improved geometry~\cite{kollmannsberger2019}. A fitted source factor also needs to be interpreted in relation to energy and observation definitions. Scanning calorimetry shows that energy coupling and penetration vary with power, speed and beam size across the conduction--keyhole transition~\cite{ye2019}; IN625 measurements distinguish unreflected laser power from energy delivered to the substrate~\cite{lane2020absorption}. In the AMMT benchmark, imaging length and metallographic width/depth are obtained through different procedures and sampled track populations~\cite{lane2020}. A length-calibrated source factor consequently represents the combination of source, model and observation mapping used in that fit, rather than a separately identified intrinsic absorption.

% I5: Existing diagnostics establish a narrower present question.
Existing assessments already examine several aspects of this interpretive problem. Comparisons of Gaussian conduction predictions with thermal imaging and metallography show that useful surface trends can coexist with depth bias, under property and phase-change assumptions that differ from a nonlinear enthalpy model~\cite{soares2026}. Analytical sensitivity studies quantify the response of geometry and temperature to processing and material inputs~\cite{yang2024}. Inverting measured bare-track depth provides a further diagnostic: it tests whether a conduction model would require physically infeasible absorption to explain the observation~\cite{hong2026}. These approaches assess different outputs and assumptions, so their results cannot be transferred directly to a particular calibrated temperature field. The present question is narrower: when high-temperature material assumptions change in a length-calibrated model, does the resulting length change represent movement of the thermal tail, a change in the spatial freezing interval, or a change in the time material spends traversing that interval?

% I6: The specific property-data boundary and storage distinction.
The property inputs make this question concrete. The IN625 conductivity and sensible specific-heat tables used here end at 982 and 1093~$^\circ$C, respectively, below the adopted solidus of 1290~$^\circ$C~\cite{specialmetals625}. The supplier describes conductivity as measured and specific heat as calculated typical alloy data; their continuation into the melting range is a modeling choice, not a measurement of the benchmark coupon's liquid properties. Conductivity determines diffusive heat redistribution, whereas sensible heat capacity contributes to enthalpy storage. In an enthalpy formulation, latent heat adds storage within the adopted phase interval~\cite{vanelsen2007}. Changing sensible heat capacity while retaining that latent contribution therefore differs from changing the entire phase-change storage law. Independent continuation changes allow the two property responses and their interaction to be examined within the same calibrated model.

% I7: Define the spatial-to-material-time relation before relying on it.
Interpreting those responses also requires distinguishing position, spatial separation and material time. Behind the translating laser, the centerline temperature crosses the adopted liquidus at 1350~$^\circ$C and solidus at 1290~$^\circ$C during cooling. The distance between these rear crossings defines the spatial phase span. Moving both crossings by similar distances can shift the thermal tail and change total pool length while leaving that span nearly unchanged. For a quasi-steady field translating at scan speed $v$, stationary substrate material passes through a span $\ell_m$ in time $\ell_m/v$. The corresponding interval-mean cooling rate follows from the temperature interval divided by this passage time. These are derived diagnostics of the same temperature field, referring to stationary-material passage rather than a liquid-parcel trajectory. Spatial geometry alone therefore does not specify the material cooling time when scan speed changes.

% I8: The actual comparison, data roles and bounded contribution.
The present study examines these distinctions for three AMMT bare-plate IN625 conditions using a nonlinear moving-frame conduction model with sensible and latent enthalpy storage. Condition B surface length alone defines the calibration objective; width/depth are excluded from fitting, and conditions A and C retain the fitted source factor. All reference targets were available during development, making the comparisons retrospective and nonblind. Model length is a two-solidus-crossing surface proxy, whereas width and depth describe the maximum-temperature fusion envelope during passage. The reference imaging length and metallographic cross-section means are separate condition-level observations~\cite{lane2020,nist2018}. At B, a fixed-source four-corner comparison independently replaces linear continuation by last-value holding for conductivity and sensible heat capacity while retaining the latent law. Rear solidus/liquidus positions then separate tail displacement from phase-span change, and the span and scan speed give material passage time and interval-mean cooling rate. This design identifies how property assumptions alter specific features of the modeled thermal history after one geometric calibration, while keeping the remaining shape discrepancy and the need for independent thermal validation visible.
